Let @@M0@@ be a free graded @@M1@@-Module of finite type. By 1.2, we can find a graded isomorphism:
\sgaStableTarget{sga6.e06.eq.1-2-2}@@M2@@

We thus obtain a decomposition of @@M3@@ as a direct sum of free graded @@M4@@-Modules of finite type that are homogeneous, namely:
@@M5@@
where @@M6@@ is free homogeneous of degree @@M7@@ and of rank @@M8@@.

\sgaStableTarget{sga6.e06.prop.1-2}\textbf{Proposition 1.2.3.} \uline{The integers} @@M9@@ \uline{occurring in decomposition} (1.2.2) \uline{do not depend on the chosen decomposition}.

Suppose given two decompositions, hence a graded isomorphism:
@@M10@@
The proof proceeds by induction on the number of nonzero free homogeneous direct summands occurring in the decomposition on the left. If this is @@M11@@, then @@M12@@, and the proposition is evident. Suppose it proved for @@M13@@, and let us prove it for @@M14@@. The given isomorphism, being compatible with the gradings, gives an isomorphism on the homogeneous components of the two sides, and in particular an isomorphism of the homogeneous components of degree @@M15@@, which shows that for @@M16@@, we have @@M17@@, and on the homogeneous components of degree @@M18@@, which shows that they have the same rank, and consequently that @@M19@@. The given isomorphism therefore induces an isomorphism of the free homogeneous direct summands of lowest degree in the two decompositions, and hence passes to the quotient by these direct summands; we then apply the induction hypothesis to the quotient Module, whence the result.

